\section{Experimental Setup}
\label{sec:setup}

\subsection{Benchmark construction}

\paragraph{Source dataset.} We build on ReproFlake~\cite{zhang2026reproflake},
a dataset of 1{,}115 machine-reproducible flaky tests (819 ID, 126 OD, 125
NIO, 36 TD, 9 uncategorized) from 63 JIRA issues and 1{,}052 commits of the
iDoFT dataset~\cite{parry2022survey}. Each test ships a project snapshot at
the flaky commit, the developer fix as a patch (\texttt{Fixed.patch}), a
Maven dependency cache, and the \emph{official reproduction mechanism} of its
category: manual sleep injection at a timing-sensitive line (TD, shipped as
	a \texttt{FlakyCodeChange.patch} reproduction patch), polluter--victim
	ordered runs (OD), the
iDFlakies same-JVM double-run plugin (NIO), and NonDex seed perturbation (ID).
The fix patch is our diagnosis ground truth; the official mechanisms validate
our subset and seed the warm-start ablation.

\paragraph{Subset selection.} We select 116 tests: \emph{all} 36 TD tests
(the scarcest, most timing-sensitive category), plus 30 OD, 25 NIO, and 25 ID
tests chosen by a greedy, download-efficient procedure with per-zip caps. The
subset deliberately balances categories (rather than mirroring ReproFlake's
ID-dominated distribution) to make per-category analysis meaningful. The
manifest records, per test: category, project, module, Java version, polluter
(OD), NonDex seed and rounds (ID), flaky commit, and fix commit, spanning 46
project snapshots (17.4\,GB).

\paragraph{Validation gate.} Every test must first \emph{pass} a gate that
runs its category's official mechanism in our environment and observes the
target test failing: TD by applying the dataset's sleep-injection patch (\texttt{FlakyCodeChange.patch}), recompiling the touched module when the
injection sits in production code, then running up to three times; OD up to 10
polluter--victim runs in the recorded order via the dataset's customized
Surefire build (\texttt{runOrder=testorder}, shipped in each project's
dependency cache); NIO up to 10 same-JVM double
runs via the iDFlakies testrunner plugin with the dataset's detector
configuration; ID one NonDex run with the dataset's recorded seed (the
mechanism is deterministic given the seed, so one run suffices). A test
fails the gate if the mechanism never reproduces within its cap, or if the
project fails to build offline (45-minute cap). Gate results are reported in
Table~\ref{tab:benchmark}; the gate rate is the fraction of \emph{built} tests
whose official mechanism reproduced the failure (build failures are reported
separately in the \emph{Built} column). Eighteen tests across seven snapshots
fail to build offline and are excluded: ten OD tests in a single
\texttt{wildfly-naming} snapshot (a dependency missing from the shipped cache),
five TD tests in Hadoop/HDFS/YARN snapshots (the Hadoop build requires a
\texttt{protoc} binary that is absent from the offline environment), and three
ID tests in a \texttt{fastjson2} snapshot (a test-compilation error under the
recorded JDK). The gate dataset is the ground-truth file each of our
reproduction mechanisms is checked against.

\begin{table}[t]
\caption{Benchmark composition and validation-gate results. Gate runs use
each category's official ReproFlake mechanism (TD: sleep-injection patch;
OD/NIO: 10 ordered/double runs; ID: one NonDex run with the dataset's seed).}
\label{tab:benchmark}
\small
\begin{tabular}{@{}lrrrrr@{}}
\toprule
\textbf{Category} & \textbf{Selected} & \textbf{Projects} & \textbf{Built} &
\textbf{Gate passed} & \textbf{Gate rate}\\
\midrule
Order-dependent (OD) & 30 & 6 & 20 & 20 & 100.0\%\\
Time-dependent (TD) & 36 & 26 & 31 & 31 & 100.0\%\\
Non-idempotent (NIO) & 25 & 5 & 25 & 25 & 100.0\%\\
Implementation-dep.\ (ID) & 25 & 9 & 22 & 16 & 72.7\%\\
\midrule
\textbf{Total} & 116 & 46 & 98 & 92 & 93.9\%\\
\bottomrule
\end{tabular}
\end{table}

\subsection{Models}

Table~\ref{tab:models} lists the six models. All are 2026 releases (the
oldest, DeepSeek-V4.1, was released after January 2026), and all are accessed
exclusively through vendor APIs — no model is deployed locally. Every LLM
call in every experiment uses temperature $0.2$ and at most 2{,}000 output
tokens, except the two DeepSeek models, which get 4{,}000: both are reasoning
models whose visible reply shares the output budget with hidden reasoning
tokens, and at 2{,}000 tokens the reply came back empty in early pilots. The
reasoning-effort knob, where offered, is set to the lowest level for the same
reason. DeepSeek-V4.1-Flash is the main model (full 92-test
gate-passed set, all ablations, variance); the other five run on a stratified 50-test
subset that preserves the category proportions of the benchmark. DeepSeek is
called directly at its off-peak list prices; OpenAI, Gemini, and Grok models
are called through a commercial OpenAI-compatible API relay, whose per-model
prices (which we use for cost accounting) are the prices in the table; for
Grok-4.7 we report the relay's \emph{stable}-channel price.

\begin{table}[t]
\caption{The six 2026-era models used, with the prices we paid (USD per 1M
tokens; relay prices for OpenAI/Gemini/Grok, off-peak list for DeepSeek).
Access is API-only; all experiments fix temperature $0.2$.}
\label{tab:models}
\small
\begin{tabular}{@{}lllrrl@{}}
\toprule
\textbf{Key} & \textbf{Model} & \textbf{Provider} & \textbf{In \$} &
\textbf{Out \$} & \textbf{Role}\\
\midrule
ds-flash & DeepSeek-V4.1-Flash & DeepSeek & 0.15 & 0.60 & main\\
ds-pro & DeepSeek-V4-Pro & DeepSeek & 0.66 & 1.98 & 50-test subset\\
gemini38 & Gemini-3.8-Flash & Google & 0.11 & 0.56 & 50-test subset\\
grok47 & Grok-4.7 & xAI & 0.22 & 0.66 & 50-test subset\\
sol & GPT-5.6 & OpenAI & 0.40 & 2.00 & 50-test subset\\
astra & GPT-6-Astra & OpenAI & 1.00 & 5.00 & 50-test subset\\
\bottomrule
\end{tabular}
\end{table}

\subsection{Baselines}

All methods — FlakeAgent and the five baselines — run on the same gate-passed
tests, in the same sandboxes, with the same 8-execution budget
(Table~\ref{tab:baselines}); the agent additionally gets up to 6 LLM calls.
The baselines are deliberately \emph{log-free}, matching the cold-start
setting: log-based reproducers such as TDRepro~\cite{chen2024tdrepro} cannot
run at all without the failing log, so they are inapplicable to this setting
(we revisit this in Section~\ref{sec:discussion}).

\begin{table}[t]
\caption{The five baselines. All share the 8-execution budget with
FlakeAgent; SingleShotLLM additionally gets one LLM call.}
\label{tab:baselines}
\small
\begin{tabular}{@{}llp{7.6cm}@{}}
\toprule
\textbf{Baseline} & \textbf{Executions} & \textbf{Description}\\
\midrule
NaiveRerun & 8 & isolated reruns of the target test (T1), stop on failure\\
RandomStrategy & 8 & uniform random (strategy, params) over the same sanitized space, stop on failure\\
OrderShuffle & 8 & full module suite in random order (T4), fresh seeds, stop on failure\\
SleepScan & 8 & deterministic sleep scan: T6 with 100/500/2000/5000\,ms $\times$2 into the first test-class method\\
SingleShotLLM & 8 & one LLM call proposes a single strategy; it is executed up to 8 times with fixed parameters (no loop, no observations fed back)\\
\bottomrule
\end{tabular}
\end{table}

\textbf{NaiveRerun} extends the agent's Step-0 reruns to the full budget;
it represents ``click re-run until it fails.'' \textbf{RandomStrategy}
isolates the value of \emph{selecting} strategies: it samples uniformly from
the identical strategy space with parameters drawn from the same sanitized
ranges, without any reasoning. \textbf{OrderShuffle} is the standard
iDFlakies-style whole-suite randomization applied to our budget.
\textbf{SleepScan} is a deterministic FlakeRake-style timing perturbation
scan. \textbf{SingleShotLLM} is the no-loop LLM pipeline: one planning call
(identical prompt, no prior observations), then the proposed strategy is
executed up to 8 times with fixed parameters — it measures how much of the
agent's value comes from the closed loop versus the model's prior alone.
RandomStrategy and OrderShuffle are seeded and repeated 3 times over a
30-test variance subset; FlakeAgent likewise runs 3 seeds on that subset.

\subsection{FlakeAgent configurations and ablations}

The main configuration is Section~\ref{sec:approach} verbatim: 3 Step-0
reruns, $\le$5 planning rounds, 1 diagnosis call on success, $\le$8
executions, 45-minute cap. Four ablations isolate design elements:
\textbf{(i) $-reflection$} — the loop is cut to a single planning round
after Step 0 (at most 2 LLM calls), i.e. no observation-feedback iteration;
\textbf{(ii) $-space4$} — the strategy space shrinks to
\{T1, T2, T5, T6\}, removing order permutation/randomization (T3/T4), stress
and environment strategies (T7--T9), parallelism (T10), network blocking
(T11), and NonDex (T12); \textbf{(iii) $-source$} — the test source excerpt
is withheld from every planning context (diagnosis still sees source);
\textbf{(iv) $Warm$} — the agent is seeded with the failure observation
produced by one invocation of the category's official mechanism and proceeds
directly to diagnosis, isolating the reproduction loop and approximating the
operating point of log-based tools.

\subsection{Metrics}

\textbf{Reproduction}: a test is \emph{reproduced} if the target test fails
at least once within budget; we report overall and per-category rates, plus
\emph{reliability} (mean $\pm$ std over seeds) on the variance subset.
\textbf{ATR} counts executions until first failure (medians and ECDFs).
\textbf{Strategy attribution} aggregates each reproduction's triggering
strategy into a category $\times$ strategy heat map. \textbf{Diagnosis} (only
for reproduced tests): top-1 \emph{category accuracy} (exact, plus a coarse
order-vs-non-order bin), \emph{evidence precision}, and \emph{fix-hit}
(diagnoses citing $\ge$1 line the developer changed). \textbf{Cost}: tokens
and USD per test and per reproduction at list prices, plus wall-clock time,
from provider-reported usage metadata, never estimates.

\subsection{Environment}

Experiments run on a Linux machine (WSL2, Ubuntu 22.04; 8 cores, 11\,GB RAM
allocated) with Maven 3.9.16 and JDK 8/11/17 installed; each project uses its
recorded Java version and builds offline against the dependency cache shipped
with its dataset zip. All 46 zips (17.4\,GB) are pre-extracted; per-zip
preparation (offline build) is cached. The LLM client reaches DeepSeek
directly and the other providers through a local HTTP proxy; no model runs
locally. Every execution, baseline or agent, goes through the same
observation probe, and every LLM call is logged with provider-reported token
usage — the numbers in Section~\ref{sec:results} are computed from these logs
by the scripts in our artifact, not by hand.
